% =====================================================================
% 第 1 章  引言
% ---------------------------------------------------------------------
% 编译说明同前：ctexart 下用 \section，ctexbook/report 下改为 \chapter。
% =====================================================================

\section{引言}

中国人民银行每季度发布一份《中国货币政策执行报告》。这份文件既复盘上一季度的货币信贷
与金融市场运行，也交代下一阶段的政策取向，篇幅通常在五十页上下。对市场而言，它是少数
几份能够同时提供数据与判断的官方文本；对研究者而言，它提供了连续二十余年、结构相对
稳定的书面沟通样本。一个自然的问题是：报告里的话，市场听进去了吗？

已有的回答大多集中在价格方向上。姜富伟等（2021）以四类宽基指数为对象，发现报告文本
情绪改善伴随发布后一日的正收益，并进一步指出，真正引起反应的是反映政策指引的部分，
而非回顾宏观经济基本面的部分。此后的一系列研究把这一思路延伸到不同场景：许魁（2023）
考察分析师报告在传导中的作用，王辉等（2023）关注预期引导与银行风险，蒋海等（2025）
提取前瞻性情绪解释流动性囤积，周阳、綦岚鑫（2024）构建领域词典检验股票、债券与外汇
市场的反应。这些工作共同确立了一个判断：报告文本含有市场尚未定价的信息。

但仍有三处可以推进。第一，结论多停留在市场总量，行业维度只作为稳健性出现，缺少对
“文本影响的是哪一类资产、哪一阶矩”的区分。第二，分段方式不尽相同，有的按句子内容
分类，有的按章节切分，口径差异使结论难以直接比较。第三，样本多止于 2019 年前后，
2020 年之后的非常规货币政策阶段较少被纳入，效应是否稳定缺乏检验。

推进这些问题会碰到一个技术障碍。报告级情绪每期只有一个取值，在行业面板中加入季度
固定效应后，它与时间虚拟变量完全共线，系数无法识别。这不是样本长度的问题，而是变量
层级与模型设定的必然结果。要在有时间固定效应的框架下识别情绪，只能依赖情绪与行业
特征的交互；而要在时间序列上识别，则必须把观测限制在信息到达的附近。

本文选择后者，把识别建立在报告发布日的事件窗上。具体而言，以 2001 年第一季度至
2026 年第二季度共 102 期报告为样本，用中文金融情感词典与情感单元法分别度量政策指引段
与宏观回顾段的情绪；以 16 个连续可得的申万一级行业构造等权市场组合，计算发布后一日、
五日以及发布前后的累计收益与日内振幅；在此基础上检验情绪对收益方向与波动的影响，
并按 2014 年前后分组考察效应的稳定性。

结果可以概括为三句。政策指引段的情绪对发布后的收益方向没有稳健影响，系数接近于零；
同一情绪显著降低发布日的波动，控制事前波动后依然成立；把样本按 2014 年切开后，
波动效应集中在前期，方向效应出现在后期且为负，但后者对个别年份较为敏感，本文仅将其
作为提示性证据。

相对于既有文献，本文的边际贡献有三。其一，把被解释变量从收益扩展到波动，并用日内
振幅与绝对收益两个代理交叉验证，说明文本的作用首先体现在收敛分歧而非推动价格。
其二，按章节结构切分指引段与回顾段，在保留姜富伟等（2021）核心判断的同时，给出了
一种可复现的替代口径。其三，引入制度断点并配合断点年份敏感性、替代词典与置换检验，
对结论的稳健性作出较为完整的交代。

需要说明的是，本文没有估计“哪些行业对政策指引更敏感”。受数据条件限制，行业特征
（估值、杠杆、利率敏感度）的覆盖期偏短，无法支撑稳定的交互估计，这一扩展留待后续
研究。相应地，本文的结论限定在市场层面，行业数据的作用是构造市场组合与提供稳健性
对照。

全文余下部分安排如下。第 2 章回顾央行沟通、文本情绪度量与行业传导三支文献，并提出
待检验的假说。第 3 章说明数据来源、文本处理与变量构造。第 4 章讨论识别策略与模型
设定。第 5 章报告实证结果与稳健性检验。第 6 章讨论机制、政策含义与局限，并给出结论。
